% Complete copied proof fragment; source and changes are in BUILD_MANIFEST.json.
\section{Finished calculations on the retained physical and arithmetic
formulas}\label{TranscriptStage:proof10:finished-calculations-on-the-retained-physical-and-arithmetic-formulas}

These calculations use the objects and constants displayed in A0656 and
A0706. They close bounded algebraic omissions. They do not identify the
modular surface, the physical wedge, or the fluid/gravity expansion with
the later split-base cohomology by assumption. The source coordinates
below are LF coordinates in the pinned transcript segment.

\subsection{1. The modular scattering coefficient retains every zero
multiplicity}\label{TranscriptStage:proof10:the-modular-scattering-coefficient-retains-every-zero-multiplicity}

A0656, lines 2964--3096, displays

\[
\mathsf S(k)=\pi^{-2ik}
\frac{\Gamma(\frac12+ik)}{\Gamma(\frac12-ik)}
\frac{\zeta(1+2ik)}{\zeta(1-2ik)},
\qquad E_* = \frac{\hbar^2}{2mR^2}>0,
\qquad \mathcal E(k)=E_*(k^2+\tfrac14).
\]

Take any nontrivial zeta zero in the specified lower half-plane, with
its actual multiplicity:

\[
\rho=a-i\gamma,\qquad 0<a<1,\quad\gamma>0,
\qquad q=\operatorname{ord}_{s=\rho}\zeta(s)\ge1.
\]

Solving the denominator equation exactly gives

\[
k_\rho=\frac\gamma2-\frac{i(1-a)}2,
\quad 1-2ik_\rho=\rho,
\quad 1+2ik_\rho=2-\rho,
\quad \tfrac12-ik_\rho=\tfrac\rho2,
\quad \tfrac12+ik_\rho=1-\tfrac\rho2.
\]

Both Gamma arguments have positive real part. Gamma is finite and
nonzero there. Also \(\Re(2-\rho)=2-a>1\), where the absolutely
convergent Euler product is nonzero. Consequently the analytic numerator

\[
N(k)=\pi^{-2ik}
\frac{\Gamma(\frac12+ik)}{\Gamma(\frac12-ik)}\zeta(1+2ik)
\]

is nonzero at \(k_\rho\). Taylor expansion at the actual zero gives

\[
\zeta(1-2ik)
=\frac{\zeta^{(q)}(\rho)}{q!}(-2i)^q(k-k_\rho)^q
+O((k-k_\rho)^{q+1}).
\]

Thus the pole order is exactly \(q\), and its leading Laurent
coefficient is exactly

\[
a_{-q}=
\frac{q!N(k_\rho)}{(-2i)^q\zeta^{(q)}(\rho)}\ne0.
\]

The map to energy has derivative \(2E_*k_\rho\ne0\), so it is locally
invertible at this pole. The energy-plane pole has the same order and
leading coefficient

\[
b_{-q}=a_{-q}(2E_*k_\rho)^q.
\]

Direct squaring, without removing either real correction, gives

\[
\mathcal E_\rho
=E_*\left[\frac{\gamma^2}{4}-\frac{(1-a)^2}{4}+\frac14\right]
-i\frac{E_*\gamma(1-a)}2.
\]

Let \(\sum_{j=1}^{q}b_{-j}(\mathcal E-\mathcal E_\rho)^{-j}\) be the
complete principal part. The positively oriented local contour
contribution with the explicitly specified convention
\((2\pi i)^{-1}\oint e^{-i\mathcal Et/\hbar}\mathsf S(k(\mathcal E))\,d\mathcal E\)
equals

\[
e^{-i\mathcal E_\rho t/\hbar}
\sum_{j=1}^{q}\frac{b_{-j}}{(j-1)!}
\left(-\frac{it}{\hbar}\right)^{j-1}.
\]

Indeed, expand the exponential in powers of
\(\mathcal E-\mathcal E_\rho\); its coefficient of power \(j-1\) is the
displayed factor, which is exactly the residue paired with the
order-\(j\) pole. A lower-half-plane contour used in a global inverse
transform may carry a different overall orientation; the convention
above fixes this sign and makes no unproved contour deformation. Every
multiple zero contributes the complete polynomial of degree \(q-1\),
with nonzero top coefficient, multiplying the decaying exponential. Its
exponential amplitude factor is

\[
\exp\!\left[-\frac{E_*\gamma(1-a)t}{2\hbar}\right].
\]

The source's simple-pole width \(\Gamma_\rho=E_*\gamma(1-a)\) is
therefore retained, and the previously undisplayed multiplicities are
now explicit. This calculation concerns the principal part of the
continued scattering coefficient. It does not assert a complete
time-domain expansion for every wavepacket, or convert poles of a
continued coefficient into nonreal eigenvalues of a self-adjoint
Hamiltonian. Those claims require a specified inversion contour, packet
class, and bounds on the remaining integral.

\subsection{2. Channel redistribution is exact and need not be
positive}\label{TranscriptStage:proof10:channel-redistribution-is-exact-and-need-not-be-positive}

A0656, lines 3104--3137, uses a Hilbert space \(\mathcal H\), a unitary
\(U:\mathcal H\to\mathcal H\), and an orthogonal projection
\(P=P^*=P^2\). For every \(\psi\in\mathcal H\),

\[
\langle\psi,(U^*PU-P)\psi\rangle
=\langle U\psi,PU\psi\rangle-\langle\psi,P\psi\rangle
=\|PU\psi\|^2-\|P\psi\|^2.
\]

Moreover,

\[
\tfrac12U^*[2P-I,U]
=\tfrac12U^*((2P-I)U-U(2P-I))
=U^*PU-P.
\]

No domain issue occurs: both operators are bounded. Unitarity does not
imply positivity of this operator. In \(\mathbb C^2\) with its usual
inner product, take

\[
P=\begin{pmatrix}1&0\\0&0\end{pmatrix},
\qquad U=\begin{pmatrix}0&1\\1&0\end{pmatrix}.
\]

Then

\[
U^*PU-P=\begin{pmatrix}-1&0\\0&1\end{pmatrix}.
\]

The two coordinate unit vectors have expectations \(-1\) and \(+1\).
This retains the precise useful interpretation of a change in the
selected observable. Any use of this identity to prove an arithmetic
positivity statement still requires the actual arithmetic subspace and
pairing; the identity by itself supplies no such sign.

\subsection{3. Equal marginal histories leave a phase family even after
cross-energy
measurement}\label{TranscriptStage:proof10:equal-marginal-histories-leave-a-phase-family-even-after-cross-energy-measurement}

A0706, lines 3835--3901, displays a thermofield vector and the
partial-trace map. Work on the exact countable energy eigenbasis
appearing there. Put

\[
p_n=\frac{e^{-\beta E_n}}{Z(\beta)},\qquad
\rho_\beta=\sum_n p_n|n\rangle\langle n|,
\qquad \sum_n p_n=1.
\]

For every real phase sequence \(\theta=(\theta_n)\), the norm-convergent
vector

\[
|\Psi_\theta\rangle=\sum_n\sqrt{p_n}e^{i\theta_n}|n\rangle_L|n\rangle_R
\]

has norm one because its summands are mutually orthogonal. In the
partial trace of \(|\Psi_\theta\rangle\langle\Psi_\theta|\), the right
basis inner product is \(\langle m|n\rangle=\delta_{mn}\); hence both
marginal density operators are exactly \(\rho_\beta\). The product
density \(\rho_\beta\otimes\rho_\beta\) has the same marginals. This
proves the source's many-to-one statement.

When \(\sum_n p_nE_n^2<\infty\), diagonal summation gives

\[
\langle H_LH_R\rangle_{\Psi_\theta}=\sum_n p_nE_n^2,
\qquad
\langle\delta H_L\delta H_R\rangle_{\Psi_\theta}
=\sum_n p_nE_n^2-\left(\sum_n p_nE_n\right)^2.
\]

This value is independent of all phases. It is zero in the product state
after centering because the expectation factors. Thus the particular
cross-energy observable in A0706 distinguishes the displayed product
state from a nonzero-variance thermofield vector, but it does not
recover the phase pattern inside the entire pure-state family.

An explicit bounded cross observable that does retain this missing phase
is

\[
O_{mn}=(|m\rangle\langle n|)_L\otimes(|m\rangle\langle n|)_R,
\qquad m\ne n.
\]

Direct action on the series gives

\[
\langle\Psi_\theta|O_{mn}|\Psi_\theta\rangle
=\sqrt{p_mp_n}\,e^{i(\theta_n-\theta_m)}.
\]

Its Hermitian real and imaginary parts measure respectively the cosine
and sine of this relative phase, with the exact factor
\(\sqrt{p_mp_n}\). This is a concrete continuation beyond the source's
single cross-energy observable.

For completeness, partial trace after arbitrary local unitary driving
obeys

\[
\operatorname{Tr}_R[(U_L\otimes U_R)\rho(U_L^*\otimes U_R^*)]
=U_L(\operatorname{Tr}_R\rho)U_L^*.
\]

To prove it, pair both sides with any bounded \(A_L\); cyclicity of the
trace replaces \((A_L\otimes I)\) by \((U_L^*A_LU_L\otimes I)\), and the
defining property of partial trace gives equality of all such
expectations. The same argument applies on the right. Therefore equal
initial marginals give equal histories for the same specified local
unitary drives. This is an obstruction to reconstructing the joint state
from those histories alone. It does not obstruct a reconstruction that
includes the displayed cross observables, a joint evolution law, or an
independently specified joint state.

\subsection{4. The arithmetic occupation map, modular clock, and
trace-class
boundary}\label{TranscriptStage:proof10:the-arithmetic-occupation-map-modular-clock-and-trace-class-boundary}

A0706, lines 3903--3980, gives the BC Hamiltonian

\[
\mathsf H|n\rangle=E_*\log n\,|n\rangle,
\qquad E_*>0.
\]

Let \(\mathfrak P\) denote the primes. The bosonic Fock occupation basis
over \(\ell^2(\mathfrak P)\) is indexed by sequences
\(m=(m_p)_{p\in\mathfrak P}\) of nonnegative integers of finite total
sum. Such a sequence has finite support. Unique factorization is a
bijection

\[
n\longleftrightarrow(v_p(n))_{p\in\mathfrak P},
\qquad (m_p)_p\longmapsto\prod_pp^{m_p}.
\]

Sending one orthonormal basis to the other gives a surjective isometry
\(\mathcal U\), hence a unitary. On its full diagonal-operator domain,

\[
\mathcal U\mathsf H\mathcal U^*=\sum_pE_*\log p\,N_p.
\]

The equality of domains is explicit: a Fock coefficient sequence \(c_m\)
belongs to this domain exactly when

\[
\sum_m\left(\sum_pE_*m_p\log p\right)^2|c_m|^2<\infty,
\]

which is the image of \(\sum_n(E_*\log n)^2|c_n|^2<\infty\). Thus this
is an operator equality, including its domain, rather than a formal
equality of symbols.

Write \(\sigma=\beta E_*\). The positive trace is

\[
\operatorname{Tr}(e^{-\beta\mathsf H})=\sum_{n\ge1}n^{-\sigma}.
\]

It is finite exactly when \(\sigma>1\), by comparison with
\(\int_1^\infty x^{-\sigma}dx\); for \(\sigma\le0\) the terms do not
even tend to zero. For any finite prime set \(F\), summing the
independent geometric occupation series gives

\[
\sum_{\operatorname{supp}(m)\subset F}\prod_{p\in F}p^{-\sigma m_p}
=\prod_{p\in F}(1-p^{-\sigma})^{-1}\quad(\sigma>0).
\]

Increasing \(F\) through all primes and using nonnegative monotone sums
recovers \(\sum_n n^{-\sigma}\). Hence the full Gibbs density and
doubled vector exist in the displayed Hilbert-space representation
exactly for \(\beta E_*>1\). In that range,

\[
|\mathrm{TFD}_\beta\rangle
=\zeta(\beta E_*)^{-1/2}\sum_{n\ge1}n^{-\beta E_*/2}|n\rangle_L|n\rangle_R,
\]

and its prime factors are exactly those displayed in A0706. Every energy
moment used in the preceding calculation is finite because
\(\sum_n n^{-\sigma}(\log n)^j<\infty\) for every fixed nonnegative
integer \(j\) and \(\sigma>1\), as follows from the same integral
comparison after \(x=e^y\).

The frequency dictionary is \(\omega_p=E_*\log p/\hbar\). Substitution
into the source's thermal-coefficient match gives

\[
\beta E_*\log p=\frac{2\pi c\omega_p}{a}
\quad\Longleftrightarrow\quad
\beta=\frac{2\pi c}{\hbar a}.
\]

No prime-dependent choice remains. Combining this equality with the
actual trace-class domain yields the precise domain of this global
vector correspondence:

\[
0<a<\frac{2\pi cE_*}{\hbar}.
\]

At the upper endpoint, \(\beta E_*=1\), the vector normalization
diverges through the harmonic series, although every individual prime's
geometric factor is finite. This is an omitted global boundary in the
retained response. It is a limitation of this trace-class arithmetic
representation, not a theorem that an arbitrary Rindler observer lacks a
KMS state.

The modular sign and proper clock also agree exactly with the source
convention. On the thermofield representation the modular operator for
the left algebra is

\[
\Delta=\rho_\beta\otimes\rho_\beta^{-1}.
\]

For a bounded left observable, the convention used in A0706 gives

\[
\Delta^{-is}(A\otimes I)\Delta^{is}
=(e^{i\beta s\mathsf H}Ae^{-i\beta s\mathsf H})\otimes I.
\]

The scalar partition factors cancel. Physical Heisenberg time is
\(\alpha_t(A)=e^{it\mathsf H/\hbar}Ae^{-it\mathsf H/\hbar}\), so
\(t=\beta\hbar s=2\pi cs/a\), exactly the source relation
\(s=a\tau/(2\pi c)\). This proves a thermal representation and clock
correspondence. It does not construct a spatial observable net, a
relativistic stress tensor, or a cooling generator.

\subsection{5. The force pullback has a real terminal-control
issue}\label{TranscriptStage:proof10:the-force-pullback-has-a-real-terminal-control-issue}

A0706, lines 3660--3697, constructs a one-form through the actual
preterminal fluid flow:

\[
A_t=-(\Phi_t^{-1})^*\int_0^t\Phi_s^*f_s\,ds,
\qquad h_{00}=-2v^jA_j.
\]

Pulling back by \(\Phi_t\), differentiating, and using the chain rule
for a transported one-form gives

\[
\Phi_t^*[(\partial_t+\mathcal L_v)A_t]=-\Phi_t^*f_t,
\quad
\partial_tA_i+v^j\partial_jA_i+A_j\partial_iv^j=-f_i.
\]

Therefore the boundary-force expression is exactly

\[
\tfrac12\partial_i h_{00}-\partial_tA_i
+v^j(\partial_iA_j-\partial_jA_i)
=-\partial_tA_i-v^j\partial_jA_i-A_j\partial_iv^j=f_i.
\]

In Lagrangian coordinates the full derivative factor is

\[
A_t(\Phi_t(x))
=-D\Phi_t(x)^{-T}
\int_0^tD\Phi_s(x)^Tf_s(\Phi_s(x))\,ds.
\]

This identifies exactly what must be computed on the source's actual
flow, rather than silently inferring terminal regularity from \(f\).

The general implication from terminal-flat forcing to smooth
coefficients is false even among smooth preterminal forced
incompressible NS solutions. The following explicit comparison tests
that implication without replacing the source's finite-energy field.
Retain arbitrary viscosity \(\nu>0\), source time \(0\le t<1\), and put

\[
u=1-t,\qquad g(t)=e^{-1/u^2},
\quad v(t,x)=\left(-\frac{x_1}{u},\frac{x_2}{u},0\right),
\quad p(t,x)=g(t)x_1-\frac{x_2^2}{u^2},
\quad f=g(t)\,dx^1.
\]

The force extends smoothly by zero to \(t=1\), with all terminal time
derivatives zero: every derivative is a polynomial in \(u^{-1}\) times
\(e^{-1/u^2}\), which tends to zero. The velocity has divergence
\(-1/u+1/u=0\) and zero Laplacian. Its first material-acceleration
component is zero and its second is \(2x_2/u^2\); adding
\(\nabla p=(g(t),-2x_2/u^2,0)\) gives exactly \(f\). Thus

\[
\partial_tv+v\cdot\nabla v-\nu\Delta v+\nabla p=(g(t),0,0)
\]

for every retained \(\nu>0\).

The flow and pulled-back force are

\[
\Phi_t(x)=(ux_1,u^{-1}x_2,x_3),
\qquad \Phi_s^*f_s=(1-s)g(s)\,dx^1.
\]

Writing \(I(t)=\int_0^t(1-s)g(s)\,ds\), the prescribed construction
yields

\[
A_t=-\frac{I(t)}u\,dx^1,
\qquad h_{00}=-\frac{2I(t)x_1}{u^2}.
\]

Here \(I(t)\) tends to the strictly positive finite number
\(\int_0^1(1-s)e^{-1/(1-s)^2}ds\), so these coefficients do not extend
smoothly to the endpoint. This example is not compactly supported and
does not have finite total kinetic energy; it therefore does not settle
the actual NS manuscript's terminal behavior. It proves exactly that the
displayed preterminal force-map identity alone cannot make that
inference. The strongest continuation remains evaluating its full flow
derivatives and the gravitational corrections on the actual retained NS
profile.
